\section*{ID6998: What to Look for in Simulation Outputs: N(t): N(t) (1)}
\paragraph{Metadata.}
PiID: 6744604774649 | RTEID: RTE:P27174.7881:T0.7934 | UUID: 1a96d8a8-7c83-5751-adda-0c7823b505e1

\paragraph{Canonical Statement.}
- If \textbackslash{}(N(t)\textbackslash{}) settles at a finite value and \textbackslash{}(N\_e(t)\textbackslash{}) is steady, you’re in continuous steady-state — compare thermalization vs loss rates. - If \textbackslash{}(N(t)\textbackslash{}) shows exponential growth, you may be crossing a stimulated-emission threshold (could be lasing or runaway gain) — check whether that’s physical (saturation not included). - If \textbackslash{}(N\_e\textbackslash{}) saturates near \textbackslash{}(N\_d/2\textbackslash{}) with positive gain, that implies inversion; for photon BEC we generally aim for strong thermalization without inversion. - Sweep `tau\_c` (Q) and `pump\_rate` to map threshold contours where stimulated term becomes comparable to cavity loss.

\paragraph{Canonical Equation.}
\[
N(t)
\]

\paragraph{Dictionary.}
Relation: N(t)

\paragraph{Encyclopedia.}
What to Look for in Simulation Outputs: N(t): N(t) (1) is a symbolic expression, written N(t). It introduces a symbol or relation used elsewhere in the framework. It is built from N. Within the Trawin Zero Point registry it serves as one element of the framework's mathematical narrative.
